\subsection{Local stochastic expansion}

Write $\mathbb G_n=\sqrt n(\mathbb P_n-P_0)$ and
$\boldsymbol{\psi}^*(\mathbf{x})=\boldsymbol{\psi}(\mathbf{x};\boldsymbol{\theta}^*)$.  Although $m_{\boldsymbol{\theta}}(\mathbf{x})$ is continuous in
$\boldsymbol{\theta}$, it is not differentiable when two component scores tie.  The
probability of observations lying in a shrinking tube around an active face
therefore controls the empirical remainder.

For $j<l$, define the activity margin
\begin{equation}
 \Delta_{jl}(\mathbf{x};\boldsymbol{\theta})
 =\min\{q_j(\mathbf{x};\boldsymbol{\theta}),q_l(\mathbf{x};\boldsymbol{\theta})\}
  -\max_{m\notin\{j,l\}}q_m(\mathbf{x};\boldsymbol{\theta}),
 \label{eq:activity-margin}
\end{equation}
where the maximum over an empty set is $-\infty$.

\begin{assumption}[Local stochastic regularity]
\label{ass:local-stochastic}
There exist a neighbourhood $U$ of $\boldsymbol{\theta}^*$, constants
$C,t_0>0$, and a measurable envelope $L$ with
$P_0L^{2+\eta}<\infty$ for some $\eta>0$ such
that:
\begin{enumerate}
 \item for all $\boldsymbol{\theta},\boldsymbol{\vartheta}\in U$ and every $j$,
 \[
  |q_j(\mathbf{X};\boldsymbol{\theta})-q_j(\mathbf{X};\boldsymbol{\vartheta})|
  +\|\nabla_{\boldsymbol{\theta}}q_j(\mathbf{X};\boldsymbol{\theta})-\nabla_{\boldsymbol{\theta}}q_j(\mathbf{X};\boldsymbol{\vartheta})\|
  \leq L(\mathbf{X})\|\boldsymbol{\theta}-\boldsymbol{\vartheta}\|
 \]
 almost surely, with the analogous local bound for $\nabla_{\boldsymbol{\theta}}^2q_j$;
 \item for every $j<l$ and $0<t<t_0$,
 \[
  P_0\{|g_{jl}(\mathbf{X};\boldsymbol{\theta}^*)|\leq tL(\mathbf{X}),\
          \Delta_{jl}(\mathbf{X};\boldsymbol{\theta}^*)\geq-tL(\mathbf{X})\}\leq Ct;
 \]
 \item simultaneous $t$-neighbourhoods of two distinct active contrasts have
 probability $O(t^2)$; and
 \item $\mathbf{A}=-\nabla_{\boldsymbol{\theta}}^2M(\boldsymbol{\theta}^*)$ is positive definite.
\end{enumerate}
\end{assumption}

The second condition is a local margin condition.  It follows under the
regular-face assumptions when the distribution of an active contrast has a
bounded density at zero, with suitable integrability for $L$.  The third
condition rules out first-order probability mass at multiple-face junctions.

\begin{lemma}[Quadratic mean expansion of the max criterion]
\label{lem:l2-expansion}
Under Assumptions~\ref{ass:boundary-geometry} and
\ref{ass:local-stochastic}, as $\mathbf{u}\to0$,
\begin{equation}
 \left\|m_{\boldsymbol{\theta}^*+\mathbf{u}}-m_{\boldsymbol{\theta}^*}-\mathbf{u}^\mathsf T\boldsymbol{\psi}^*\right\|_{P_0,2}
 =O(\|\mathbf{u}\|^{3/2})+O(\|\mathbf{u}\|^2).
 \label{eq:l2-expansion}
\end{equation}
The bound holds uniformly over $\mathbf{u}/\|\mathbf{u}\|$.
\end{lemma}

\begin{proof}
Put $r_{\mathbf u}=m_{\boldsymbol\theta^*+\mathbf u}-
m_{\boldsymbol\theta^*}-\mathbf u^\mathsf T\boldsymbol\psi^*$ and let
$T_{\mathbf u}$ be the union of the active tubes in
Assumption~\ref{ass:local-stochastic}(2), with $t=C_0\|\mathbf u\|$.
On $T_{\mathbf u}^c$ the winner is unchanged, and Taylor's theorem gives
$|r_{\mathbf u}|\leq C L\|\mathbf u\|^2$. On $T_{\mathbf u}$, the Lipschitz
bound gives $|r_{\mathbf u}|\leq C L\|\mathbf u\|$. After the standard
truncation implied by $P_0L^{2+\eta}<\infty$, the margin condition yields
\begin{align*}
 P_0r_{\mathbf u}^2
 &\leq C\|\mathbf u\|^4P_0L^2
   +C\|\mathbf u\|^2P_0(L^2\ind\{T_{\mathbf u}\})\\
 &=O(\|\mathbf u\|^4)+O(\|\mathbf u\|^3).
\end{align*}
Taking square roots proves \eqref{eq:l2-expansion}. Multiple-tube intersections
are $O(\|\mathbf u\|^2)$ by part (3), so they do not alter the bound. A full
bracketing construction and the precise truncation argument appear in
Supplementary Lemma 1.3.
\end{proof}

\begin{theorem}[Local empirical-process expansion]
\label{thm:local-empirical-expansion}
Under Assumptions~\ref{ass:boundary-geometry} and
\ref{ass:local-stochastic}, for every finite $R$,
\begin{multline}
 \sup_{\|\mathbf{u}\|\leq R}\left|
 n\left[(M_n-M)(\boldsymbol{\theta}^*+\mathbf{u}/\sqrt n)
       -(M_n-M)(\boldsymbol{\theta}^*)\right]
 -\mathbf{u}^\mathsf T\mathbb G_n\boldsymbol{\psi}^*
 \right|\\
 \xrightarrow{p}0.
 \label{eq:local-empirical-expansion}
\end{multline}
Consequently,
\begin{equation}
 n\{M_n(\boldsymbol{\theta}^*+\mathbf{u}/\sqrt n)-M_n(\boldsymbol{\theta}^*)\}
 =\mathbf{u}^\mathsf T\mathbb G_n\boldsymbol{\psi}^*
  -\tfrac12\mathbf{u}^\mathsf TAu+o_p(1),
 \label{eq:local-asymptotic-quadratic}
\end{equation}
uniformly on compact subsets of $\mathbb R^d$.
\end{theorem}

\begin{proof}
Write $r_{n,\mathbf u}=r_{\mathbf u/\sqrt n}$. Directly,
\begin{align*}
 &n(\mathbb P_n-P_0)\{m_{\boldsymbol\theta^*+\mathbf u/\sqrt n}
                       -m_{\boldsymbol\theta^*}\}
 -\mathbf u^\mathsf T\mathbb G_n\boldsymbol\psi^*\\
 &\hspace{35mm}=\sqrt n\,\mathbb G_n r_{n,\mathbf u}.
\end{align*}
Lemma~\ref{lem:l2-expansion} gives
$\sup_{\|\mathbf u\|\leq R}\|r_{n,\mathbf u}\|_{P_0,2}=O(n^{-3/4})$.
Supplementary Lemma 1.3 bounds the bracketing entropy by
$C d\log(C/\epsilon)$; the bracketing maximal inequality therefore gives
$\sup_{\|\mathbf u\|\leq R}|\sqrt n\mathbb G_nr_{n,\mathbf u}|=o_p(1)$.
This proves \eqref{eq:local-empirical-expansion}. Finally, Taylor expansion of
$M$ gives, uniformly on compact sets,
\[
 n\{M(\boldsymbol\theta^*+\mathbf u/\sqrt n)-M(\boldsymbol\theta^*)\}
 =\tfrac12\mathbf u^\mathsf T\mathbf H\mathbf u+o(1)
 =-\tfrac12\mathbf u^\mathsf T\mathbf A\mathbf u+o(1),
\]
which yields \eqref{eq:local-asymptotic-quadratic}
\citep{vandervaart1998,van1996weak}.
\end{proof}

\subsection{Rate, asymptotic linearity, and normality}

Define
\begin{equation}
 \mathbf{J}=P_0\{\boldsymbol{\psi}^*(\mathbf{X})\boldsymbol{\psi}^*(\mathbf{X})^\mathsf T\}.
 \label{eq:score-variance}
\end{equation}
At the interior maximizer, Theorem~\ref{thm:population-curvature} gives
$P_0\boldsymbol{\psi}^*=0$.

\begin{theorem}[Root-$n$ rate]
\label{thm:root-n-rate}
Suppose the conditions of Theorems~\ref{thm:cml-consistency} and
\ref{thm:local-empirical-expansion} hold.  If the approximate-maximization
error in \eqref{eq:approximate-maximizer} satisfies $nr_n=o_p(1)$, then
\begin{equation}
 \sqrt n(\widehat{\boldsymbol{\theta}}_n-\boldsymbol{\theta}^*)=O_p(1).
 \label{eq:root-n-rate}
\end{equation}
\end{theorem}

\begin{proof}
Let $\lambda_0=\lambda_{\min}(\mathbf A)>0$. For sufficiently small
$\delta$, the population expansion implies
$M(\boldsymbol\theta^*+\mathbf v)-M(\boldsymbol\theta^*)
\leq-(\lambda_0/4)\|\mathbf v\|^2$ whenever $\|\mathbf v\|\leq\delta$.
On the dyadic shell
$2^jC<\|\mathbf u\|\leq2^{j+1}C$, the local empirical term is
$O_p(2^jC)$ whereas the negative drift is at most
$-(\lambda_0/4)2^{2j}C^2$. Hence, uniformly over shells inside the
$\delta$-neighbourhood,
\[
 P\!\left(\sup_{\|\mathbf u\|>C}
 n\{M_n(\boldsymbol\theta^*+\mathbf u/\sqrt n)-M_n(\boldsymbol\theta^*)\}
 \geq-nr_n\right)\longrightarrow0
\]
as first $n\to\infty$ and then $C\to\infty$. Consistency excludes the
complement of that neighbourhood, and $nr_n=o_p(1)$ does not change the shell
comparison. Thus $\|\sqrt n(\widehat{\boldsymbol\theta}_n-
\boldsymbol\theta^*)\|=O_p(1)$.
\end{proof}

\begin{theorem}[Asymptotic linearity and normality]
\label{thm:asymptotic-normality}
Under the conditions of Theorem~\ref{thm:root-n-rate},
\begin{equation}
 \sqrt n(\widehat{\boldsymbol{\theta}}_n-\boldsymbol{\theta}^*)
 =\mathbf{A}^{-1}\mathbb G_n\boldsymbol{\psi}^*+o_p(1).
 \label{eq:asymptotic-linearity}
\end{equation}
If $\mathbf{J}$ is finite, then
\begin{equation}
 \sqrt n(\widehat{\boldsymbol{\theta}}_n-\boldsymbol{\theta}^*)
 \rightsquigarrow N_d(0,\mathbf{V}),
 \qquad \mathbf{V}=\mathbf{A}^{-1}\mathbf{J}\mathbf{A}^{-\mathsf T}.
 \label{eq:asymptotic-normality}
\end{equation}
This result remains centered on the CML functional $\boldsymbol{\theta}^*$ when the
ordinary mixture model is misspecified.
\end{theorem}

\begin{proof}
By Theorem~\ref{thm:root-n-rate}, the local parameter
$\widehat{\mathbf{u}}_n=\sqrt n(\widehat{\boldsymbol{\theta}}_n-\boldsymbol{\theta}^*)$ is tight. Put
$\mathbf Z_n=\mathbb G_n\boldsymbol\psi^*$. Completing the square in the local
criterion gives
\[
 \mathbf u^\mathsf T\mathbf Z_n-	frac12\mathbf u^\mathsf T\mathbf A\mathbf u
 =\tfrac12\mathbf Z_n^\mathsf T\mathbf A^{-1}\mathbf Z_n
 -\tfrac12(\mathbf u-\mathbf A^{-1}\mathbf Z_n)^\mathsf T
 \mathbf A(\mathbf u-\mathbf A^{-1}\mathbf Z_n).
\]
Uniformity on compact sets, tightness, and $nr_n=o_p(1)$ force
$\widehat{\mathbf u}_n-\mathbf A^{-1}\mathbf Z_n=o_p(1)$, proving
\eqref{eq:asymptotic-linearity}. Since
$\mathbf Z_n\rightsquigarrow N_d(\mathbf0,\mathbf J)$ by the multivariate
central limit theorem, Slutsky's theorem gives covariance
$\mathbf A^{-1}\mathbf J\mathbf A^{-\mathsf T}$.
\end{proof}

\begin{proposition}[Transfer to a local algorithmic solution]
\label{prop:algorithmic-transfer}
Let $\widetilde{\boldsymbol{\theta}}_n\to_p\boldsymbol{\theta}^*$ be an interior solution produced by an
algorithm.  Suppose ties occur with probability zero and
\begin{equation}
 \mathbb P_n\boldsymbol{\psi}(\mathbf{X};\widetilde{\boldsymbol{\theta}}_n)=o_p(n^{-1/2}),
 \label{eq:algorithmic-score-error}
\end{equation}
and suppose the local score process is stochastically differentiable at
$\boldsymbol{\theta}^*$ with derivative $-\mathbf{A}$.  Then
\begin{equation}
 \sqrt n(\widetilde{\boldsymbol{\theta}}_n-\boldsymbol{\theta}^*)
 =\mathbf{A}^{-1}\mathbb G_n\boldsymbol{\psi}^*+o_p(1).
 \label{eq:algorithmic-linearity}
\end{equation}
Thus a consistent stationary solution has the same first-order distribution
as the local CML maximizer when its optimization error is smaller than the
statistical score error.
\end{proposition}

\begin{proof}
Stochastic differentiability gives
\[
 \mathbb P_n\boldsymbol{\psi}(\mathbf{X};\widetilde{\boldsymbol{\theta}}_n)
 =\mathbb P_n\boldsymbol{\psi}^*-\mathbf{A}(\widetilde{\boldsymbol{\theta}}_n-\boldsymbol{\theta}^*)
  +o_p(n^{-1/2}+\|\widetilde{\boldsymbol{\theta}}_n-\boldsymbol{\theta}^*\|).
\]
Use $P_0\boldsymbol{\psi}^*=0$, multiply by $\sqrt n$, and solve the resulting linear
equation.  Positive definiteness of $\mathbf{A}$ and consistency absorb the remainder.
\end{proof}

\begin{theorem}[Validity of the nonparametric bootstrap]
\label{thm:bootstrap-validity}
Let $\mathbb P_n^{\#}$ be the empirical measure of a size-$n$ sample drawn
with replacement from $\mathbb P_n$, and let
$\widehat{\boldsymbol{\theta}}_n^{\#}$ be a bootstrap approximate maximizer
initialized in the same consistent basin as
$\widehat{\boldsymbol{\theta}}_n$.  Suppose, conditionally on the data, its
optimization error is $o_p(n^{-1})$, and suppose the localized classes in
Theorem~\ref{thm:local-empirical-expansion} are bootstrap-Donsker and
$\mathbf{J}$ is positive definite.  Then
\begin{equation}
 \sup_{C\in\mathcal C}
 \left|
 P^{\#}\!\left\{
 \sqrt n(\widehat{\boldsymbol{\theta}}_n^{\#}
          -\widehat{\boldsymbol{\theta}}_n)\in C
 \right\}
 -P\{N_d(\mathbf{0},\mathbf{V})\in C\}
 \right|\xrightarrow{p}0,
 \label{eq:bootstrap-validity}
\end{equation}
where $\mathcal C$ is the collection of convex Borel subsets of
$\mathbb R^d$.  Thus the ordinary nonparametric bootstrap consistently
estimates the first-order distribution, provided each bootstrap fit reaches
the relevant local solution.
\end{theorem}

\begin{proof}
Let $\mathbb G_n^\#=\sqrt n(\mathbb P_n^\#-\mathbb P_n)$. Bootstrap
equicontinuity of the localized remainder class gives, conditionally in
probability and uniformly for bounded $\mathbf u$,
\begin{align*}
 &n\{M_n^\#(\widehat{\boldsymbol\theta}_n+\mathbf u/\sqrt n)
       -M_n^\#(\widehat{\boldsymbol\theta}_n)\}\\
 &=\mathbf u^\mathsf T\mathbb G_n^\#
   \boldsymbol\psi(\cdot;\widehat{\boldsymbol\theta}_n)
   -\tfrac12\mathbf u^\mathsf T\widehat{\mathbf A}_n\mathbf u
   +o_{p^\#}(1).
\end{align*}
Here $\widehat{\mathbf A}_n\to_p\mathbf A$, and the conditional central limit
theorem gives
$\mathbb G_n^\#\boldsymbol\psi(\cdot;\widehat{\boldsymbol\theta}_n)
\rightsquigarrow_{p}N_d(\mathbf0,\mathbf J)$. Completing the square, as in the
proof of Theorem~\ref{thm:asymptotic-normality}, and using the conditional
$o_p(n^{-1})$ optimization error yields
\[
 \sqrt n(\widehat{\boldsymbol\theta}_n^\#-\widehat{\boldsymbol\theta}_n)
 =\widehat{\mathbf A}_n^{-1}\mathbb G_n^\#
 \boldsymbol\psi(\cdot;\widehat{\boldsymbol\theta}_n)+o_{p^\#}(1).
\]
The conditional Slutsky theorem gives the Gaussian limit with covariance
$\mathbf V$. Since that law is nondegenerate, its boundary assigns zero mass
to convex-set continuity boundaries, yielding the displayed uniform metric
\citep{van1996weak}.
\end{proof}
